\section{Cross-Layer ECC Evaluation Methodology}

\begin{figure*}[t]
    \centering
    \includegraphics[width=0.98\textwidth]{figure01_methodology.pdf}
    \caption{Identity-preserving evaluation workflow. Insufficient evidence exits the comparison or remains explicitly unavailable; it is never imputed.}
    \label{fig:method}
\end{figure*}

\subsection{Correctness and Identity Qualification}

Each candidate binds a code identity to a specific RTL architecture and interface. Only candidates with sufficient functional evidence enter comparison. Qualification records the payload/codeword relation, correction and detection guarantee, latency, initiation interval (II), and the correspondence between the tested implementation and the physically characterized RTL. For the two conventional SECDED candidates, exhaustive formal equivalence after a two-cycle temporal alignment establishes identical extended-Hamming (72,64) input/output semantics. Thus, their physical difference is a microarchitectural treatment rather than a different code.

This step prevents nominal labels such as ``SECDED'' or ``BCH'' from silently pooling measurements from unrelated implementations. It also prevents a physically successful but functionally unqualified design from becoming a comparison point.

The comparison contract fixes the transaction boundary as well as the code. Encoder and decoder channels are independent, parallel interfaces with no backpressure, and each can accept one valid transaction per cycle. Every implementation carries the same 64 useful payload bits, while the realized 72- or 78-bit codeword width is retained as an implementation cost. Latency is measured from accepted request to corresponding response; II describes successive acceptance. Consequently, a three-cycle design with II=1 is neither mislabeled as lower-throughput nor treated as equal-latency to a one-cycle design.

Qualification and characterization are deliberately separated. Correctness evidence determines admission; it cannot be rescued by favorable PPA. Conversely, a candidate that lacks physical evidence is not assigned a synthetic worst-case point. This distinction is important for Pareto analysis: an unavailable coordinate is epistemically different from a measured but unattractive coordinate.

\subsection{Reliability Semantics}

We separate guarantees from observations. The conventional and Hsiao SECDED implementations correct all single-coordinate errors and detect all double-coordinate errors. The shortened BCH (78,64,$t=2$) implementation corrects all one- and two-coordinate errors in its proven universe.

For context, we also exhaustively enumerate weight-3 masks over canonical codeword coordinates and classify decoded outcomes as silent data corruption (SDC) or detected uncorrectable error (DUE). These are finite-universe observations beyond the guaranteed correction radius. They are not FIT, SER, operational probabilities, or field-weighted rates. In particular, BCH spans $\binom{78}{3}$ masks while SECDED spans $\binom{72}{3}$; fractions must not be treated as samples from one common physical fault distribution.

\subsection{Common Physical Characterization}

Accepted RTL enters one pinned OpenROAD-flow-scripts (ORFS) environment. Technology, corner, target clock, I/O delays, output load, utilization policy, placement density, seed, and worker count are common. The area metric is final standard-cell instance area; timing and detailed wirelength are post-route. Achieved frequency retains negative slack and is calculated as
\begin{equation}
 f_{\max}=\frac{1000}{T_{\mathrm{target}}-\mathrm{WNS}}\ \mathrm{MHz}.
 \label{eq:fmax}
\end{equation}
Routing completion does not imply target feasibility: a zero-DRC result with negative setup slack remains a routed but timing-infeasible point.

\subsection{Equal-Work Power and Energy}

Power uses post-route OpenSTA with deterministic VCD activity. Each timing-feasible implementation receives the same no-error workload of 100,000 accepted 64-bit transactions at II=1 and a 10 ns clock. Both traces contain six reset, 100,000 valid-work, and four drain cycles; four drain cycles exceed the three-cycle pipeline boundary. Input annotation covers the same 139 interface pins before internal propagation. Energy per useful operation is
\begin{equation}
 E_{\mathrm{op}}=P\,T_{\mathrm{trace}}/N_{\mathrm{useful}}.
 \label{eq:energy}
\end{equation}
Equal duration and operation count make the power and energy percentage changes identical. This is a steady-stream comparison, not a claim of lower request latency.

The activity experiment also preserves the implementation boundary. Identical input-pin annotation is propagated into each design, so reported totals include the distinct internal switching induced by its RTL and placed-and-routed netlist. We do not decompose the measured difference into register, clock-tree, or combinational contributions; doing so would require an additional experiment. The result therefore supports an empirical equal-work energy contrast, not a causal claim that pipelining generally reduces ECC energy.

\subsection{Decision Rule}

We avoid a weighted scalar score because its weights could conceal categorical differences in correction guarantee and timing feasibility. Comparisons instead use exact pairwise enumeration. Within one guarantee tier, design $a$ dominates $b$ only if it is no worse in every admitted objective and strictly better in at least one. The timing-feasible SECDED objective set minimizes area, latency, and achievable energy per operation while maximizing Fmax. Detailed wirelength is reported as a routing-cost diagnostic and included in sensitivity checks.

Correction guarantees form separate strata: W1 correction/W2 detection and W1/W2 correction are not collapsed into a numeric benefit. A stronger guarantee is therefore not declared globally superior when it is physically worse, and a physically dominating SECDED design does not erase BCH's stronger guarantee. Timing-infeasible designs retain routed area, wirelength, slack, and achieved Fmax, but they cannot contribute target-clock energy. This rule makes the boundary between measured evidence and design preference explicit.
